\begin{solapartat}
\punts{0,2} Segons la llei de gravitació universal, el mòdul de la força sobre un satèl·lit de massa $m$ que orbita al voltant de la Terra amb un radi orbital $r$ s'expressa com a:
\[ F = G\,\frac{m M_\mathrm{T}}{r^2} \]

\punts{0,2} I la segona llei de Newton estableix que: $\vec{F} = m\vec{a}$

\punts{0,2} D'altra banda, considerant que el satèl·lit descriu un moviment circular uniforme al voltant de la Terra, la seva acceleració centrípeta és: $a = \dfrac{v^2}{r}$

\punts{0,25} Com que sobre el satèl·lit només hi actua la força de la gravetat:
\[ G\,\frac{m M_\mathrm{T}}{r^2} = m\,\frac{v^2}{r} \;\Rightarrow\; v = \sqrt{G\,\frac{M_\mathrm{T}}{r}} \]

\punts{0,2} La velocitat orbital del nanosatèl·lit és:
\[ v = \sqrt{G\,\frac{M_\mathrm{T}}{r}} = \sqrt{6,67\times 10^{-11}\,\frac{5,98\times 10^{24}}{6,37\times 10^{6} + 500\times 10^{3}}} = 7620\ \mathrm{m/s} \]

\punts{0,2} I el període és:
\[ T = \frac{2\pi r}{v} = \frac{2\pi(6,37\times 10^{6} + 500\times 10^{3})}{7620} = 5660\ \mathrm{s} \]
\end{solapartat}

\begin{solapartat}
\punts{0,4} Com que el camp gravitatori és conservatiu, l'energia mecànica del satèl·lit a la superfície de la Terra és igual a l'energia mecànica a l'òrbita.

De l'apartat anterior sabem la velocitat orbital:
\[ v = \sqrt{G\,\frac{M_\mathrm{T}}{r}} \]

\punts{0,1} Per tant, l'energia cinètica per a una òrbita de radi $r$ és:
\[ E_c(r) = \frac{1}{2}mv^2 = \frac{1}{2}G\,\frac{M_\mathrm{T} m}{r} \]

\punts{0,2} L'energia potencial a l'òrbita és:
\[ E_p(r) = -G\,\frac{M_\mathrm{T} m}{r} \]

\punts{0,2} I l'energia mecànica a l'òrbita és:
\[ E_m(r) = E_c(r) + E_p(r) = -\frac{1}{2}G\,\frac{M_\mathrm{T} m}{r} \]

\punts{0,1} A la superfície de la Terra, l'energia potencial és:
\[ E_p(R_\mathrm{T}) = -G\,\frac{M_\mathrm{T} m}{R_\mathrm{T}} \]

I si igualem les energies mecàniques obtenim:
\[ \frac{1}{2}mv^2 - G\,\frac{M_\mathrm{T} m}{R_\mathrm{T}} = -\frac{1}{2}G\,\frac{M_\mathrm{T} m}{r} \;\Rightarrow\; \frac{1}{2}v^2 = G M_\mathrm{T}\left(\frac{1}{R_\mathrm{T}} - \frac{1}{2r}\right) \]

\punts{0,1} I, finalment, si aïllem la velocitat obtenim:
\[ v = \sqrt{G M_\mathrm{T}\left(\frac{2}{R_\mathrm{T}} - \frac{1}{r}\right)} \]

\punts{0,15} Per al nanosatèl·lit:
\[ v = \sqrt{6,67\times 10^{-11}\cdot 5,98\times 10^{24}\left(\frac{2}{6,37\times 10^{6}} - \frac{1}{6,37\times 10^{6} + 500\times 10^{3}}\right)} = 8200\ \mathrm{m/s} \]
\end{solapartat}
